\documentclass{situationist-preprint}

\title{The Aesthetic Value of Ugliness}
\msid{AES-12-a}
\msfamily{AES-12}
\msdiscipline{Aesthetics}
\mstype{Principal manuscript}
\mskeywords{ugliness; negative aesthetic properties; aesthetic value; organic unity; expression; disgust}

\begin{document}
\maketitle

\begin{abstract}
Can ugly art be good? If ugliness is the contrary of beauty, and beauty is what makes art good, ugly art is bad art, and the many admired works that are ugly must be admired for something other than their aesthetic character. We argue that ugly art can be good, and that its goodness can be aesthetic. Ugliness is an aesthetic property, response-dependent in the way beauty is: an object is ugly if it is disposed to give aesthetic displeasure to qualified observers in suitable conditions. Like other aesthetic properties, it contributes to the value of a work as a whole according to its role in the work, and not according to its value in isolation. Where a work's prescribed response, such as horror at suffering or revulsion at corruption, can be secured only through features that are ugly, the ugliness is a means to the work's aesthetic achievement, and good art may deliberately exploit it. We distinguish this view from the claim that art may represent ugly things beautifully, which is true but does not answer the question.
\end{abstract}

\section{The problem}

The crucified Christ of Grünewald's Isenheim Altarpiece has a body covered with sores, twisted hands and a greenish, decomposing skin. Goya's late paintings on the walls of his house show grotesque, distorted figures; Francis Bacon's popes and portraits smear and contort the human face. These are among the most admired works of European painting, and they are ugly: not only in what they represent, but in how they look. If the aim of art is beauty, they should be failures.

The traditional answer saves the classical view by a distinction. Aristotle observed that we take pleasure in accurate images of things whose sight is painful, such as the forms of the lowest animals and of corpses (\emph{Poetics} 4). Kant held that art shows its superiority by describing beautifully things that in nature would be ugly or displeasing, with the single exception of what arouses disgust \citep[§48]{Kant1790}. On this view the Isenheim Christ is a beautiful representation of an ugly body. \citet{Rosenkranz1853} gave ugliness a systematic treatment but kept it subordinate: the ugly has a place in art as a negative moment, justified by its resolution in a higher beauty, as in the comic. More recently, \citet{Korsmeyer2011} has argued that disgust, the response Kant excluded, can be an aesthetic emotion that art legitimately arouses and that contributes to its value.

The traditional distinction is correct as far as it goes, but it does not cover the cases. Bacon's paintings are not beautiful pictures of ugly faces. Their facture is itself harsh, and the ugliness is in the painting. We argue that ugliness is an aesthetic property that good art may deliberately exploit.

\section{Framework}

\begin{definition}[Ugliness]\label{d:ugly}
An object is \emph{ugly} if and only if it is disposed to give aesthetic displeasure to qualified observers who perceive it, in suitable conditions, in its correct category.
\end{definition}

Definition~\ref{d:ugly} is the counterpart of the response-dependent account of beauty defended in AES-02. Aesthetic displeasure is the displeasure taken in an object as perceived, and not in its consequences or its usefulness. Ugliness is not the mere absence of beauty: many objects are neither beautiful nor ugly, and ugliness has a positive phenomenology of its own, from the jarring to the repellent.

\begin{principle}[Aesthetic properties]\label{p:aesth}
A property is aesthetic if it is perceived in an object, its possession supervenes on the object's non-aesthetic features together with its category, and its attribution is answerable to responses of qualified observers.
\end{principle}

\begin{principle}[Holism]\label{p:holism}
The contribution of an aesthetic property to the aesthetic value of a work depends on its role in the work as a whole, and need not have the same sign as its value in isolation.
\end{principle}

Principle~\ref{p:holism} is an instance of the principle of organic unities: the value of a whole is not the sum of the values of its parts \citep{Moore1903}. In art the principle is familiar. A dissonance is unpleasant in isolation and may be what makes a cadence; a harsh color may be what gives a composition its tension.

\begin{principle}[Prescribed response]\label{p:presc}
A work's aesthetic success includes its success in securing the responses it prescribes from a qualified audience (AES-08).
\end{principle}

\section{The argument}

\begin{lemma}\label{l:ugly}
Ugliness is an aesthetic property.
\end{lemma}

\begin{proof}
By Definition~\ref{d:ugly}, an object's ugliness is perceived in it, being a disposition to give aesthetic displeasure on perception; it is answerable to the responses of qualified observers; and, as a disposition, it supervenes on the categorical features that ground it, which are the object's non-aesthetic features as perceived in its category. By Principle~\ref{p:aesth}, it is an aesthetic property.
\end{proof}

\begin{lemma}\label{l:means}
If a work prescribes a response that can be secured, in that work, only through features that are ugly, the ugliness of those features contributes positively to its aesthetic value.
\end{lemma}

\begin{proof}
Suppose a work prescribes a response $R$ and secures it through features that are ugly, and that without them it would not secure $R$. By Principle~\ref{p:presc}, securing $R$ is part of the work's aesthetic success. The ugly features are then a condition of that success, and by Principle~\ref{p:holism}, their contribution to the value of the whole is fixed by this role, not by the displeasure they give in isolation. Their contribution is therefore positive.
\end{proof}

\begin{theorem}\label{t:main}
Ugly art can be good. Ugliness is an aesthetic property, and good art may deliberately exploit it as a means to its aesthetic achievement.
\end{theorem}

\begin{proof}
By Lemma~\ref{l:ugly}, ugliness is an aesthetic property. Many works prescribe responses, such as horror at physical suffering, revulsion at moral corruption or the disquiet of violence done to the human form, that cannot be secured by beautiful features: a beautiful rendering of suffering invites contemplation where the work prescribes horror. The Isenheim Christ was painted for a hospital whose patients suffered from a disease that ravaged the skin, and it prescribes recognition of suffering as bodily and unredeemed by grace of form. A work of this kind can secure its prescribed response only through ugliness, and by Lemma~\ref{l:means}, its ugliness contributes positively to its aesthetic value. Such a work can be good, and good on aesthetic grounds. Since the means are chosen to secure a prescribed response, the exploitation is deliberate.
\end{proof}

The traditional view is a special case. When a work represents an ugly object beautifully, its beauty serves a prescribed response that tolerates beauty, such as contemplative pity. Theorem~\ref{t:main} adds the cases the traditional view cannot handle, in which the prescribed response does not tolerate beauty and the work must itself be ugly.

\section{Objections and replies}

\begin{objection}[Contradiction]\label{o:contra}
If ugliness is a negative aesthetic property, a work that is ugly is to that extent aesthetically bad. Theorem~\ref{t:main} makes ugliness both an aesthetic defect and an aesthetic merit.
\end{objection}

\begin{reply}
It makes ugliness a property whose value in isolation is negative and whose contribution in context may be positive, which Principle~\ref{p:holism} allows and the case of dissonance shows to be ordinary. There is no contradiction in a property's being displeasing as perceived and valuable for what it does in a whole.
\end{reply}

\begin{objection}[Not really ugly]\label{o:transfig}
A great ugly work is not experienced as ugly. Its ugliness is transfigured by its context, so that what we see is a terrible beauty \citep{Rosenkranz1853}.
\end{objection}

\begin{reply}
If ugliness were transfigured into beauty, the work would no longer secure the response that required its ugliness, and Lemma~\ref{l:means} would not apply; but then the work would fail on its own terms. That we admire Bacon's paintings and continue to find them repellent shows that the displeasure persists in the admiration. What is transfigured is our judgment of the whole, not our perception of the parts.
\end{reply}

\begin{objection}[Disgust]\label{o:disgust}
Kant held that disgust cannot be represented without destroying aesthetic satisfaction, since the representation is felt as the thing itself. The extreme of ugliness is therefore excluded from art \citep[§48]{Kant1790}.
\end{objection}

\begin{reply}
Theorem~\ref{t:main} claims only that ugliness \emph{may} be exploited, not that every degree of it can be. Whether disgust can serve a prescribed response is a further question, and \citet{Korsmeyer2011} has given reasons to think it can. If Kant is right about some cases, the limit falls at the point where displeasure blocks attention to the work, and so blocks the prescribed response; this is a limit of the kind Lemma~\ref{l:means} predicts.
\end{reply}

\section{Conclusion}

Ugly art can be good. Ugliness is an aesthetic property, response-dependent as beauty is, and like other aesthetic properties it contributes to a work's value according to its role in the whole. Where a work prescribes a response that only ugly features can secure, its ugliness is part of its aesthetic achievement, and good art exploits it deliberately. The view keeps the traditional insight that art can render ugly things beautifully, and adds that art can be ugly, and good for being so.

\begin{thebibliography}{99}\small
\bibitem[Kant(1790)]{Kant1790} Kant, I. (1790). \emph{Critique of the Power of Judgment}. Trans. P. Guyer and E. Matthews. Cambridge University Press, 2000.
\bibitem[Korsmeyer(2011)]{Korsmeyer2011} Korsmeyer, C. (2011). \emph{Savoring Disgust: The Foul and the Fair in Aesthetics}. Oxford University Press.
\bibitem[Moore(1903)]{Moore1903} Moore, G.~E. (1903). \emph{Principia Ethica}. Cambridge University Press.
\bibitem[Rosenkranz(1853)]{Rosenkranz1853} Rosenkranz, K. (1853). \emph{Ästhetik des Häßlichen}. Gebrüder Bornträger.
\end{thebibliography}

\end{document}
